% Propietario conservado: /Users/ruben/Documents/New project/output/EXTREMA_MEDIA_RAZON_RECIPROCIDAD_ES_EN_20260918/ARTICULO/ES/source/libro/02_norma_ternaria.tex
\chapter{Norme algébrique ternaire et rayons optimaux de récupération}
\label{x_reciprocidad:lib02:norma-ternaria}

Le rayon de récupération des deux mémoires contient le facteur
\(\sqrt{73}\). Son origine peut être suivie sans introduire de grandeur
cible : il provient du Gram du lecteur nonadique et, dans une autre réalisation du
cycle ternaire déjà construit, c’est l’échelle d’une similitude. Son carré,
\(73\), est l’indice de l’inclusion réticulaire de rang deux. Ces
apparitions sont réunies par des opérations explicites ;
elles ne sont pas identifiées par une ressemblance décimale avec alpha. La comparaison
ultérieure avec \(10^4\alpha_A\) sera traitée comme un problème inverse,
avec ses données, ses transformations permises et ses conclusions exactes.

La matière première conserve sa production APP--TRIT--TPK et la coupure ultérieure décrite dans \cref{x_reciprocidad:lib01:registro}. Le registre entier, l'orientation des déplacements, la charge et les lecteurs sont reçus avant le calcul des distances. Le cas canonique utilise \(8+1\) et le dénominateur neuf. Le paramètre \(a>0\) introduit ci-dessous est une famille algébrique ultérieure, et non un nouveau germe ni un paramètre physique sélectionné par alpha. Les erreurs appartiennent au canal de lecture ; elles ne sont pas identifiées à des événements ou à des histoires admissibles.

\section{Opérateurs paramétriques à double mémoire et matrice de Gram}
\label{x_reciprocidad:lib02:familia-memoria}

Dans \(V=\mathbb R^{12}\), soit \((Sx)_i=x_{i+1}\), avec des indices modulo
douze. Pour \(a>0\) et \(j=3,4\), nous définissons
\begin{equation}
 T_{j,a}=\frac{aI+S^j}{a+1},\qquad
 D_{j,a}=\frac{\sqrt a(S^j-I)}{a+1},\qquad
 M_ax=(D_{3,a}x,D_{4,a}T_{3,a}x).
 \label{x_reciprocidad:lib02:lectores}
\end{equation}
La seconde mémoire est produite après la première compression. Archiver
les deux n’équivaut pas à réinjecter le premier complément dans une évolution
réversible à deux canaux. Soient
\[
 p(a)=a^2+a+1,\quad A(a)=a(a^2+1),\quad B(a)=a^2,
 \qquad\Pi=I-\frac{\mathbf1\mathbf1^*}{12},\quad P_u=I-\Pi.
\]

\begin{proposition}[Bilan et forme de différences]
\label{x_reciprocidad:lib02:gram}
Pour tout \(a>0\),
\[
 T_{j,a}^*T_{j,a}+D_{j,a}^*D_{j,a}=I,\qquad
 G_a:=M_a^*M_a=I-(T_{4,a}T_{3,a})^*(T_{4,a}T_{3,a}),
\]
et
\begin{align}
 (a+1)^4G_a={}&4ap(a)I
 -A(a)(S^3+S^4+S^8+S^9)\notag\\
 &-B(a)(S+S^5+S^7+S^{11}).
 \label{x_reciprocidad:lib02:gram-explicito}
\end{align}
De manière équivalente,
\begin{align}
 (a+1)^4\|M_ah\|^2={}&A(a)\sum_i
 \bigl((h_i-h_{i+3})^2+(h_i-h_{i+4})^2\bigr)\notag\\
 &+B(a)\sum_i
 \bigl((h_i-h_{i+1})^2+(h_i-h_{i+5})^2\bigr).
 \label{x_reciprocidad:lib02:energia-aristas}
\end{align}
\end{proposition}
\begin{proof}
L'orthogonalité de \(S\) donne
\[
 T_{j,a}^*T_{j,a}
 =\frac{(a^2+1)I+a(S^j+S^{-j})}{(a+1)^2},\qquad
 D_{j,a}^*D_{j,a}
 =\frac{a(2I-S^j-S^{-j})}{(a+1)^2}.
\]
Leur somme est \(I\) ; appliquer l’identité aux deux étapes annule le
terme intermédiaire \(T_{3,a}^*T_{3,a}\). Comme tous les opérateurs
sont des polynômes de \(S\), le produit des deux Grams individuels
a pour diagonale \((a^2+1)^2\), des termes simples de poids
\(a(a^2+1)\) et des termes croisés de poids \(a^2\). L’identité
\((a+1)^4-(a^2+1)^2=4a(a^2+a+1)\) donne
\eqref{x_reciprocidad:lib02:gram-explicito}. Développer les différences au carré donne
\eqref{x_reciprocidad:lib02:energia-aristas} ; chaque arête non orientée de chaque pas
apparaît une seule fois et le degré pondéré est \(4A+4B=4ap(a)\).
\end{proof}

Dans \(a=8\), on obtient
\begin{align}
 6561G_8={}&2336I-520(S^3+S^4+S^8+S^9)\notag\\
 &-64(S+S^5+S^7+S^{11}),
 \label{x_reciprocidad:lib02:gram-nonadico}
\end{align}
dont la première ligne est
\((2336,-64,0,-520,-520,-64,0,-64,-520,-520,0,-64)/6561\).
On a aussi
\[
 T_{4,8}T_{3,8}=\frac{64I+8S^3+8S^4+S^7}{81}.
\]
Les quatre déplacements distincts produisent la contribution diagonale
\[
 1-\frac{64^2+8^2+8^2+1}{81^2}
 =\frac{2336}{6561}
 =\frac{32(8^2+8+1)}{9^4}.
\]

\begin{theorem}[Noyau, spectre et inverse complet]
\label{x_reciprocidad:lib02:inversa-memorias}
Pour tout \(a>0\), \(\ker M_a=\mathbb R\mathbf1\). Les valeurs propres de \(G_a\), indexées par les modes douzièmes \(m\), sont
\[
\begin{array}{c|c|c}
 m&\lambda_m(G_a)&\text{multiplicité}\\\hline
 0&0&1\\
 3,9&2a/(a+1)^2&2\\
 4,8&3a/(a+1)^2&2\\
 6&4a/(a+1)^2&1\\
 1,5,7,11&a(5a^2+4a+5)/(a+1)^4&4\\
 2,10&a(7a^2+2a+7)/(a+1)^4&2
\end{array}
\]
En particulier,
\begin{equation}
 L_a^{\rm rec}=(G_a+P_u)^{-1}M_a^*,\qquad
 L_a^{\rm rec}M_a=\Pi,\qquad
 \|L_a^{\rm rec}\|=\frac{a+1}{\sqrt{2a}}.
 \label{x_reciprocidad:lib02:inversa}
\end{equation}
\end{theorem}
\begin{proof}
Si \(M_ah=0\), la première mémoire implique \(S^3h=h\), donc
\(T_{3,a}h=h\) ; la seconde implique \(S^4h=h\). Comme
\(S=S^4(S^3)^{-1}\), le vecteur est uniforme. La réciproque est
immédiate. Dans le mode \(m\), le bilan donne
\[
 \lambda_m(G_a)=1-
 \frac{[a^2+1+2a\cos(\pi m/2)]
 [a^2+1+2a\cos(2\pi m/3)]}{(a+1)^4}.
\]
Les douze paires de cosinus donnent le tableau. Par rapport à \(2a/(a+1)^2\), les différences des deux dernières valeurs sont
\[
 \frac{3a(a^2+1)}{(a+1)^4}>0,\qquad
 \frac{a(5a^2-2a+5)}{(a+1)^4}>0.
\]
Les autres valeurs positives sont des multiples plus grands. Le minimum
dans \(\mathbf1^\perp\) est donc \(2a/(a+1)^2\). L’opérateur
\(G_a+P_u\) vaut un sur l’uniforme et est positif inversible sur
son complément. La formule écrite est l’inverse des moindres carrés
et sa norme est l’inverse de la plus petite valeur singulière non nulle de \(M_a\).
\end{proof}

En \(a=1\), \(T_{3,1}=(I+S^3)/2\) est singulier, mais \eqref{x_reciprocidad:lib02:inversa} reste valable. La singularité d'une compression individuelle n'élimine pas la récupération au moyen de ses deux mémoires. En \(a=8\), la norme de l'inverse est \(9/4\).

\section{Distance minimale dans le domaine des registres entiers}
\label{x_reciprocidad:lib02:redes}

Les domaines algébriques de décodage sont
\[
 \mathcal C_0(a)=M_a\mathbb Z^{12},\qquad
 \Lambda_q=\{k\in\mathbb Z^{12}:\mathbf1^*k=q\},\qquad
 \mathcal C_q(a)=M_a\Lambda_q.
\]
La première image s’identifie au réseau \(\Pi\mathbb Z^{12}\)
ou à des classes modulo \(\mathbb Z\mathbf1\) ; la seconde conserve une
charge entière fixée. Ce sont des domaines qui contiennent les publications entières
du type indiqué. Leur utilisation n’affirme pas que chaque candidat corresponde à une
histoire HMT produite.

\begin{lemma}[Coupe minimale du graphe de mémoire]
\label{x_reciprocidad:lib02:cortes}
Toute coupe propre non vide des douze positions possède au moins quatre arêtes de pas \(1,5\) et au moins quatre de pas \(3,4\). Son poids dans \eqref{x_reciprocidad:lib02:energia-aristas} est donc au moins \(4A(a)+4B(a)=4ap(a)\) ; une coupe d'une position atteint cette borne.
\end{lemma}
\begin{proof}
Chaque pas \(1\) et \(5\) forme un cycle connexe à douze positions ;
chaque cycle franchit une coupe non triviale au moins deux fois. Pour les pas
\(3,4\), si l’ensemble n’est invariant par aucun des deux, chaque famille
apporte au moins deux arêtes. S’il est invariant par \(S^3\), il est une union
de classes modulo trois ; chacun des quatre triangles de \(S^4\)
le traverse deux fois, donnant huit arêtes. S’il est invariant par \(S^4\),
il est une union de classes modulo quatre ; chacun des trois cycles de longueur
quatre de \(S^3\) le traverse au moins deux fois, donnant au moins six arêtes.
L’invariance simultanée imposerait l’invariance par \(S\), impossible
pour un ensemble propre non vide. Tous les poids sont positifs ; la
coupe d’une seule position a exactement quatre arêtes de chaque
poids.
\end{proof}

\begin{theorem}[Distances exactes et minimiseurs à charge fixe]
\label{x_reciprocidad:lib02:distancias}
Pour tout \(a>0\) et toute charge entière \(q\),
\begin{equation}
 d_0(a)^2=\min_{h\notin\mathbb Z\mathbf1}\|M_ah\|^2
 =\frac{4ap(a)}{(a+1)^4},\qquad
 d_q(a)^2=\frac{8ap(a)}{(a+1)^4}.
 \label{x_reciprocidad:lib02:distancias-exactas}
\end{equation}
Dans le premier minimum, \(h\) parcourt \(\mathbb Z^{12}\) ; le second
est la distance minimale entre des mots distincts de \(\mathcal C_q(a)\).
Les différences qui atteignent le second minimum sont exactement
\(e_i-e_j\) avec \(j-i\equiv2,6,10\pmod{12}\).
\end{theorem}
\begin{proof}
Pour un entier \(h\) non constant, \((h_i-h_j)^2\geq|h_i-h_j|\). Chaque différence absolue se décompose en somme des coupes des ensembles de niveau \(\{i:h_i>t\}\). Il existe au moins un niveau non trivial ; d'après \cref{x_reciprocidad:lib02:cortes}, l'énergie pondérée vaut au moins \(4ap(a)\). Le vecteur \(e_i\) atteint cette énergie. Diviser par \((a+1)^4\) prouve le premier minimum sur tous les entiers, et pas seulement sur les indicatrices de sous-ensembles.

À charge fixe, la différence \(h\ne0\) est entière et de somme nulle.
Sa norme au carré est paire car \(h_i^2\equiv h_i\pmod2\). Si
\(\|h\|^2\geq4\), le minimum spectral donne
\[
 \|M_ah\|^2\geq\frac{8a}{(a+1)^2}
 >\frac{8ap(a)}{(a+1)^4},
\]
puisque \((a+1)^2-p(a)=a>0\). Il ne reste que
\(\|h\|^2=2\), c’est-à-dire \(h=e_i-e_j\). Son énergie est
\(2((G_a)_{ii}-(G_a)_{ij})\). Les éléments hors diagonale sont nuls
aux séparations \(2,6,10\), négatifs de poids \(B(a)/(a+1)^4\)
en \(1,5,7,11\), et négatifs de poids \(A(a)/(a+1)^4\) en
\(3,4,8,9\). La diagonale est \(4ap(a)/(a+1)^4\) ; cela prouve le
minimum et la classification exhaustive sur toute la demi-droite \(a>0\).
\end{proof}

\begin{theorem}[Rayons stricts de correction des mémoires]
\label{x_reciprocidad:lib02:radios-memoria}
Pour une donnée \(y=M_ak+e\), avec une erreur dans la norme conjointe des deux mémoires, les rayons uniformes optimaux sont
\begin{equation}
 r_0(a)=\frac{\sqrt{ap(a)}}{(a+1)^2},\qquad
 r_q(a)=\frac{\sqrt{2ap(a)}}{(a+1)^2}.
 \label{x_reciprocidad:lib02:radios}
\end{equation}
Si \(\|e\|<r_0(a)\), le plus proche voisin de \(\mathcal C_0(a)\) récupère exactement \(\Pi k\). Avec une charge \(q\) connue et \(\|e\|<r_q(a)\), le plus proche voisin de \(\mathcal C_q(a)\) récupère exactement \(k\).
\end{theorem}
\begin{proof}
Le réseau \(\Pi\mathbb Z^{12}\subset(1/12)\mathbb Z^{12}\) est
discret et \(M_a\) est minoré sur
\(\mathbf1^\perp\). Ses images sont discrètes et admettent un voisin
le plus proche. Si \(c'\) est un autre mot de code et \(d\) est la
distance minimale correspondante,
\(\|y-c'\|\geq d-\|e\|>\|e\|\) lorsque \(\|e\|<d/2\).
Le minimum est unique et la préimage est unique dans chaque domaine indiqué.
Le point médian de deux mots à distance minimale est à \(d/2\)
des deux : aucun décodeur ne peut distinguer uniformément les deux
entrées à partir de cette même donnée. L’optimalité se rapporte aux domaines
algébriques complets, non à un sous-ensemble supplémentaire d’histoires qui
pourrait avoir une séparation plus grande.
\end{proof}

Dans le cas canonique,
\begin{equation}
 r_0(8)=\frac{2\sqrt{146}}{81}\simeq0.298346814162829,
 \qquad
 r_q(8)=\frac{4\sqrt{73}}{81}\simeq0.421926110879878.
 \label{x_reciprocidad:lib02:radio73}
\end{equation}
Le \(4\) provient de \(\sqrt{2\cdot8}\) en prenant la moitié de la
distance entre registres de charge fixe ; \(73=8^2+8+1\) provient du
Gram ; \(81=9^2\) provient des deux normalisations nonadiques. Le même
produit \(9\times9\) compte les cellules APP, mais ces deux rôles
ne sont pas identifiés sans leur application typée. Le facteur quatre ne
prouve pas non plus à lui seul une assertion sur les directions spatiales.

Si les mémoires sont stockées sous la forme \(z=M_ak/\sqrt a\), la norme du bruit change. En \(a=8\), les rayons de cette donnée rationnelle sont \(r_0/\sqrt8=\sqrt{73}/81\) et \(r_q/\sqrt8=\sqrt{146}/81\). Utiliser les rayons de \eqref{x_reciprocidad:lib02:radio73} sans cette conversion revient à changer le problème. L'application à \(Q_K\), \(\eta_K\), \(B_K\) et à l'écran d'incidence a été prouvée dans \cref{x_reciprocidad:lib01:correccion-incidencia} ; en particulier, le seuil exact \(K_4=729\) n'empêche pas sa restitution discrète.

\section{Décodeur fini et dualité du paramètre}
\label{x_reciprocidad:lib02:computacion}

\begin{proposition}[Plancher et plafond sans tableau cible]
\label{x_reciprocidad:lib02:decodificador}
Avec une charge connue, sous le rayon \(r_q(a)\), il suffit d'examiner au plus \(2^{12}=4096\) candidats entiers. Sans charge, sous \(r_0(a)\), il suffit d'examiner au plus \(12\cdot2^{12}=49152\) candidats répartis entre douze résidus de charge. Le minimum résiduel renvoie respectivement \(k\) ou sa classe uniforme.
\end{proposition}
\begin{proof}
Avec la charge \(q\), formons
\(\widetilde k=L_a^{\rm rec}y+(q/12)\mathbf1\). Alors
\[
 \|\widetilde k-k\|
 <\|L_a^{\rm rec}\|r_q(a)
 =\frac{\sqrt{p(a)}}{a+1}<1.
\]
Chaque coordonnée entière vraie appartient à
\(\{\lfloor\widetilde k_i\rfloor,
\lceil\widetilde k_i\rceil\}\). On énumère ces produits,
on élimine ceux dont la somme est différente de \(q\) et on minimise
\(\|M_ak'-y\|\). Le candidat correct est présent et
\cref{x_reciprocidad:lib02:radios-memoria} assure son unicité. Les coordonnées déjà entières
réduisent le nombre de candidats.

Sans charge, écrivons \(\rho\in\{0,\ldots,11\}\) pour le reste
de la charge. Chaque classe modulo \(\mathbb Z\mathbf1\) a un
unique représentant dont la somme est \(\rho\), obtenu en soustrayant
\(\lfloor q/12\rfloor\mathbf1\). Pour chaque \(\rho\), formons
\(\widetilde k_\rho=L_a^{\rm rec}y+(\rho/12)\mathbf1\).
Dans le reste correct,
\[
 \|\widetilde k_\rho-k_\rho\|
 <\|L_a^{\rm rec}\|r_0(a)
 =\frac{\sqrt{p(a)}}{\sqrt2(a+1)}<1.
\]
On applique le plancher et le plafond, on conserve la somme \(\rho\) et on minimise le résidu entre les douze groupes. L'unicité de la classe de code donne une unique sortie centrée. Ces procédures utilisent le lecteur et l'intégralité ; elles ne reçoivent pas \(K\) comme tableau de comparaison.
\end{proof}

Pour \(a\) rationnel, on peut conserver une arithmétique rationnelle en écrivant
\(R_a=M_a/\sqrt a\), \(z=y/\sqrt a\). La reconstruction centrée
est
\[
 (R_a^*R_a+P_u)^{-1}R_a^*z=L_a^{\rm rec}y.
\]
Le résiduel peut être comparé au moyen de ses carrés rationnels. Ces
garanties concernent douze coordonnées : elles n’affirment ni complexité
polynomiale pour une dimension arbitraire, ni authentification cryptographique ni
gain quantique. L’exemple perturbé de
\cref{x_reciprocidad:lib01:correccion-incidencia} illustre l’algorithme canonique.

\begin{proposition}[Dualité réciproque et maximum du rayon]
\label{x_reciprocidad:lib02:dualidad-a}
La substitution \(a\mapsto1/a\) conserve le Gram, les distances,
les rayons et la norme de l’inverse. On a
\[
 D_{j,1/a}=D_{j,a},\qquad
 T_{j,1/a}=S^jT_{j,a}^*.
\]
Les rayons maximaux de cette famille ne sont atteints qu'en \(a=1\) :
\[
 r_q(a)\leq\frac{\sqrt6}{4},\qquad
 r_0(a)\leq\frac{\sqrt3}{4}.
\]
\end{proposition}
\begin{proof}
Les deux identités opératorielles résultent de la multiplication du numérateur et du
dénominateur par \(a\). Les gramiens individuels et la formule
\eqref{x_reciprocidad:lib02:gram} sont invariants sous la réciprocité. Posons
\(t=a/(a+1)^2\in(0,1/4]\). Alors
\[
 r_q(a)^2=2t(1-t),\qquad r_0(a)^2=t(1-t).
\]
Les deux expressions croissent pour \(0<t\leq1/4\), et \(t=1/4\)
équivaut à \(a=1\). La réciprocité laisse \(t\) fixe. Lorsque
\(a\) tend vers zéro ou l’infini, \(t\) et les rayons tendent vers zéro.
\end{proof}

L'égalité des matrices de Gram ne signifie pas que \(M_a=M_{1/a}\) : le second complément inclut une compression différente. Les extrémités \(a=0\) et \(a=\infty\) n'appartiennent pas au domaine du théorème. Optimiser la famille en \(a=1\), ou utiliser le représentant réciproque \(1/8\), ne change pas rétroactivement l'origine canonique \(a=8\).

\section{La norme dans le module d’augmentation ternaire}
\label{x_reciprocidad:lib02:modulo-ternario}

Le cycle \(r\mapsto4r\pmod9\), d’orbites \((1,4,7)\) et
\((2,8,5)\), fournit le module d’augmentation dont la carte et la forme sont
\[
 C=\begin{pmatrix}0&-1\\1&-1\end{pmatrix},\qquad
 G_{A_2}=\begin{pmatrix}2&-1\\-1&2\end{pmatrix},\qquad
 C^2+C+I=0,\qquad C^TG_{A_2}C=G_{A_2}.
\]
L’action et son orientation contragrédiente appartiennent aux fibres
ternaires antécédentes. Sur ce module déjà construit, on compose
\[
 Q_-(a)=aI-C,\qquad Q_+(a)=(a+1)I+C.
\]

\begin{proposition}[Similitude ternaire et indice entier]
\label{x_reciprocidad:lib02:norma-A2}
Pour \(a>0\),
\[
 Q_-Q_+=p(a)I,\qquad Q_-^{\dagger_{G_{A_2}}}=Q_+,
 \qquad Q_-^TG_{A_2}Q_-=p(a)G_{A_2},\qquad \det Q_-=p(a).
\]
En \(a=8\), les deux cartes conjuguées sont
\[
 Q_-=
 \begin{pmatrix}8&1\\-1&9\end{pmatrix},\qquad
 Q_+=\begin{pmatrix}9&-1\\1&8\end{pmatrix},\qquad Q_-Q_+=73I.
\]
Chaque quotient entier est cyclique d’ordre 73. Dans le quotient de \(Q_-\),
\(C\) agit par \(8\) ; dans celui de \(Q_+\), par
\(-9\equiv64\equiv8^{-1}\pmod{73}\).
\end{proposition}
\begin{proof}
Développer le produit et substituer \(C^2=-C-I\) donne \(p(a)I\).
L’invariance de la forme implique
\(C^{\dagger_{G_{A_2}}}=C^{-1}=C^2=-I-C\), d’où l’on obtient
l’adjoint. Composer le produit et l’adjoint donne la similitude. Le déterminant
de \(\bigl(\begin{smallmatrix}a&1\\-1&a+1\end{smallmatrix}\bigr)\) est \(p(a)\). En \(a=8\), le plus grand commun diviseur des entrées de chaque matrice est un et son déterminant est 73 ; la forme de Smith est \(\operatorname{diag}(1,73)\). La relation \(8I-C=0\) dans le premier quotient et \(9I+C=0\) dans le second donne les actions indiquées. Comme \(8^3\equiv1\pmod{73}\) et \(8\not\equiv1\), les deux actions sont d'ordre trois et d'orientations inverses.
\end{proof}

Ainsi, \(\sqrt{73}\) est une échelle de similitude dans la métrique ternaire
héritée, et non seulement un chiffre dans un rayon. Sa reconnaissance complexe
ultérieure est \(|8-\omega|^2=73\), avec
\(\omega^2+\omega+1=0\). L’unité complexe décrit la réalisation
de l’opérateur ; elle n’a pas été utilisée pour sélectionner le registre. L’opérateur
\(Q_-\) n’est pas identifié par ces égalités à \(K\), à toute
la dynamique TPK ou au conoyau du lecteur d’échelle d’action.

\section{Contribution invariante du cycle complet}
\label{x_reciprocidad:lib02:ciclo-completo}

Sur les douze positions, soit \(U=S^4\). Il y a maintenant quatre cycles de
longueur trois. Leurs projecteurs sont
\[
 P_0=\frac{I+U+U^2}{3},\qquad P=I-P_0.
\]
Le rang de \(P_0\) est quatre et celui de \(P\) huit. En particulier,
\(P_0\ne P_u\), car ce dernier ne conserve que l’uniforme global ;
on ne confond pas non plus \(P\) avec le projecteur exceptionnel \(P_3\).
Nous utiliserons maintenant \(Q_{-,a}=aI-U\) et \(Q_{+,a}=(a+1)I+U\)
pour les opérateurs du cycle complet.

\begin{theorem}[Matrice de Gram ternaire avec secteur fixe]
\label{x_reciprocidad:lib02:gram-ternario}
On a
\begin{equation}
 Q_{-,a}^*Q_{-,a}
 =(a-1)^2P_0+p(a)P=p(a)I-3aP_0.
 \label{x_reciprocidad:lib02:gram-Qmenos}
\end{equation}
Pour \(v\ne0\), soient
\[
 w(v)=\frac{\|P_0v\|^2}{\|v\|^2},\qquad
 s_{-,a}(v)=\frac{\|Q_{-,a}v\|^2}{\|v\|^2}.
\]
Alors \(s_{-,a}(v)=p(a)-3aw(v)\). Dans le cas canonique,
\begin{equation}
 Q_-^*Q_-=49P_0+73P=73I-24P_0,
 \qquad s(v)=73-24w(v)\in[49,73].
 \label{x_reciprocidad:lib02:73-24w}
\end{equation}
\end{theorem}
\begin{proof}
De \(U^3=I\) et \(U^*=U^2\), il résulte que \(P_0\) est le projecteur orthogonal sur \(\ker(U-I)\). En développant la matrice de Gram,
\[
 Q_{-,a}^*Q_{-,a}=(a^2+1)I-a(U+U^2).
\]
La relation \(U+U^2=3P_0-I\) donne \eqref{x_reciprocidad:lib02:gram-Qmenos}. L'évaluation sur \(v\) et la division par le carré de sa norme démontrent les autres identités.
\end{proof}

Le terme \(24\) est exactement \(3\cdot8\) : c’est la différence
\(p(8)-(8-1)^2\) entre la lecture d’augmentation et la lecture fixe. Il apparaît déjà
dans un seul cycle de trois positions. Pour \(m\) cycles, le rang fixe
est \(m\), la dimension est \(3m\) et
\[
 \operatorname{tr}(Q_{-,a}^*Q_{-,a})=3m(a^2+1).
\]
Dans \(m=4,a=8\), la trace est 780 et la correction totale de la trace
est \(24\cdot4=96\). L’égalité numérique \(24=2\cdot12\) n’est pas
l’origine opératorielle du coefficient.

Les deux cartes complètes distinguent également leurs déterminants. Dans un seul cycle,
\[
 \det(8I-U)=8^3-1=511=7\cdot73,
 \qquad
 \det(9I+U)=9^3+1=730=10\cdot73.
\]
Le facteur fixe est, respectivement, 7 et 10 ; le module d'augmentation apporte 73. Sur les douze positions, \(\det(8I-U)=511^4=68184176641\). On conserve l'identité
\begin{equation}
 73=8^2+8+1=9^2-9+1=\frac{9^3+1}{9+1}
 =\frac{729+1}{10}.
 \label{x_reciprocidad:lib02:completacion73}
\end{equation}
Son interprétation comme indice et facteur de déterminant précède
toute lecture de complétion ou de raffinement ultérieur.

\section{Six états reçus et conservation de la lecture complète}
\label{x_reciprocidad:lib02:seis-estados}

Nous fixons de nouveau \(a=8\) et omettons cet indice dans les lecteurs.
Les six objets comparés sont
\[
 K,\quad k=\Pi K,\quad m_0=D_3K,\quad m_1=D_4T_3K,
 \quad u=P_3k,\quad q_K=Q_KK=k-u.
\]
Aucun n’est sélectionné à partir de la cible de l’essai. Pour donner les
évaluations sans les remplacer par des décimales, écrivons
\(a_v=\|v\|^2\), \(b_v=\|P_0v\|^2\). On obtient
\begin{equation}
\begin{array}{c|c|c}
 v&a_v&b_v\\\hline
 K&4251257&10684363/3\\
 k&11789915/12&1170761/4\\
 m_0&15413392/81&16838032/243\\
 m_1&1105759232/6561&0\\
 u&1327717/4+(568896/5)\sqrt5
   &3029657/20+(145120/3)\sqrt5\\
 q_K&1951691/3-(568896/5)\sqrt5
   &2292404/15-(145120/3)\sqrt5
\end{array}
 \label{x_reciprocidad:lib02:tabla-normas}
\end{equation}
Le tableau s'obtient en appliquant \(P_0=(I+S^4+S^8)/3\), les formules de \(D_3,D_4T_3\) et le projecteur fini de caractère \eqref{x_reciprocidad:lib01:P3} à \eqref{x_reciprocidad:lib01:K}. Toutes ses entrées appartiennent à \(\mathbb Q(\sqrt5)\) ; les quatre premières sont rationnelles. En particulier, les trois premiers poids et lectures sont
\begin{align*}
 w(K)&=\frac{10684363}{12753771},&
 s(K)&=\frac{224866857}{4251257},\\
 w(k)&=\frac{3512283}{11789915},&
 s(k)&=\frac{776369003}{11789915},\\
 w(m_0)&=\frac{1052377}{2890011},&
 s(m_0)&=\frac{61904585}{963337}.
\end{align*}
Pour \(u\) et \(q_K\), les expressions exactes sont
\(w=b_v/a_v\), \(s=73-24b_v/a_v\), avec les numérateurs et
dénominateurs de \eqref{x_reciprocidad:lib02:tabla-normas}. Les évaluations ultérieures
sont
\[
\begin{array}{c|c|c}
 v&w(v)\text{, approximativement}&s(v)\text{, approximativement}\\\hline
 K&0.8377414805&52.89420446705527\\
 k&0.2979057101&65.85026295779062\\
 m_0&0.3641429046&64.26057028848679\\
 m_1&0&73\text{ exact}\\
 u&0.4428244530591254&62.37221312658099\\
 q_K&0.1127385160275117&70.29427561533972
\end{array}
\]
L’égalité de la seconde mémoire est imposée :
\(P_0D_4T_3=\sqrt8P_0(U-I)T_3/9=0\), tandis que
\(\|m_1\|^2>0\). Ce n’est pas une coïncidence d’évaluation.

\begin{proposition}[La mémoire complète ne modifie pas le poids fixe]
\label{x_reciprocidad:lib02:memoria-ternaria}
L'application
\[
 V_2v=(t(v),m_0(v),m_1(v))
 =(T_4T_3v,D_3v,D_4T_3v)
\]
est isométrique et entrelace \(P_0\) et \(Q_-=8I-U\) avec leurs
copies diagonales. Elle conserve séparément \(\|v\|^2\),
\(\|P_0v\|^2\) et \(\|Q_-v\|^2\). Les lectures
conjointes \(w\) et \(s\) ne changent donc pas.
\end{proposition}
\begin{proof}
Le bilan de deux étapes donne \(V_2^*V_2=I\). Chaque canal est un polynôme en \(S\), si bien qu'il commute avec \(U,P_0,Q_-\). Appliquer l'isométrie à \(v\), \(P_0v\) et \(Q_-v\) donne les trois égalités de normes. Après division, la lecture totale est la moyenne des lectures des canaux pondérée par les carrés de leurs normes, et non leur moyenne arithmétique non pondérée. Les canaux nuls ne contribuent pas au quotient total.
\end{proof}

Le terminal confirme matériellement cette partition. Pour \(K\),
\[
 \|t(K)\|^2=\frac{25538253193}{6561},\qquad
 \|P_0t(K)\|^2=\frac{848595371}{243},
\]
\[
 w(t(K))=\frac{22912075017}{25538253193},\qquad
 s(t(K))=\frac{1314402682681}{25538253193}
 \simeq51.46799480558349.
\]
La somme des trois normes au carré est \(4251257\), et celle des
normes fixes est
\((848595371+16838032)/243=10684363/3\). Pour \(k=\Pi K\),
\[
 \|t(k)\|^2=\frac{16367568169}{26244},\qquad
 \|P_0t(k)\|^2=\frac{217142795}{972}.
\]
Les compléments ne changent pas lors du centrage et leurs sommes récupèrent les normes
de la seconde ligne de \eqref{x_reciprocidad:lib02:tabla-normas}. Écarter le terminal
ou un complément altère la lecture partielle ; les conserver n’introduit pas
une correction cachée.

Dans une mémoire infinie complète, la même conclusion est maintenue lorsque
les canaux commutent avec \(U\) : la limite positive et sa racine
commutent également avec lui. Sans cette commutation, on utilise le transport
\(\mathcal VP_0\mathcal V^*\) dans l’image, conformément à
\cref{x_reciprocidad:lib01:isometria}, et non une copie diagonale supposée.

\section{Inversion, cartes et bilan bilatéral}
\label{x_reciprocidad:lib02:bidireccional}

\begin{proposition}[Ce que conserve une inversion ou un changement de repère]
\label{x_reciprocidad:lib02:inversion-carta}
Remplacer \(U\) par \(U^{-1}\) échange chaque \(Q_{\pm,a}\) avec son adjoint et conserve \(P_0\), leurs matrices de Gram, valeurs singulières, déterminants et distances induites. Pour tout \(v\ne0\), cela conserve \(w(v)\) et \(s_{-,a}(v)\). Ce n'est pas la même chose qu'inverser \(Q_{-,a}\), dont l'inverse, pour \(a\ne1\), est
\begin{equation}
 Q_{-,a}^{-1}
 =\frac1{a-1}P_0+\frac1{p(a)}Q_{+,a}P.
 \label{x_reciprocidad:lib02:inversa-Q}
\end{equation}
Pour une carte inversible \(R\), en transportant
\(U'=RUR^{-1}\), \(v'=Rv\) et
\(g'=(R^{-1})^*R^{-1}\), on conserve ces rapports de normes dans
la métrique \(g'\).
\end{proposition}
\begin{proof}
Comme \(U^*=U^{-1}=U^2\), l’inversion produit les adjoints.
Les opérateurs sont des polynômes normaux en \(U\) ; leurs Grams dépendent
de \(U+U^2\), qui ne change pas. L’inverse écrit se vérifie par
secteurs : \(Q_{-,a}\) agit par \(a-1\) dans le secteur fixe et le produit
avec \(Q_{+,a}\) vaut \(p(a)I\) dans l’augmentation. Pour la carte,
\(Q'v'=RQv\) et \(R^*g'R=I\) ; le numérateur et le dénominateur retrouvent
exactement leurs valeurs originales. L’adjoint relativement à \(g'\)
est également le transport de l’adjoint original.
\end{proof}

Le renversement de l’ordre des douze positions satisfait
\(RSR^{-1}=S^{-1}\) et conserve \(P_0\). Les translations \(S^j\),
y compris l’antipodale \(S^6\), commutent avec \(U\) ; une négation
globale n’altère pas non plus les normes. Aucune de ces opérations ne convertit
à elle seule 73 en une lecture distincte.

Avec la métrique euclidienne fixe, un repère orthogonal appliqué seulement à l'état conserve \(s_{-,a}(v)\) pour tous les \(v\) si et seulement si \(R^*P_0R=P_0\), d'après \eqref{x_reciprocidad:lib02:gram-Qmenos} et la polarisation. Le groupe de ces transformations est \(O(\operatorname{im}P_0)\times O(\operatorname{im}P)
\cong O(4)\times O(8)\), plus grand que le commutant de \(U\). En revanche, un filtre d'état \(c_0P_0+c_1P\), lorsque sa sortie est non nulle, produit
\begin{equation}
 w'=
 \frac{|c_0|^2w}{|c_0|^2w+|c_1|^2(1-w)}.
 \label{x_reciprocidad:lib02:filtro-pesos}
\end{equation}
Il peut changer la lecture. Choisir ces coefficients à partir d’une cible est
un ajustement déclaré du filtre, non un changement de carte de la même
lecture.

\begin{theorem}[Bilan bilatéral et récupération explicite]
\label{x_reciprocidad:lib02:balance-bilateral}
Pour tout \(a>0\), dans le cycle complet, on a
\begin{align}
 Q_{+,a}-Q_{-,a}^*&=3P_0,&
 Q_{+,a}Q_{-,a}&=p(a)I-3P_0,\notag\\
 Q_{+,a}^*Q_{+,a}&=p(a)I+3(a+1)P_0,
 \label{x_reciprocidad:lib02:dos-cartas}
\end{align}
et
\begin{equation}
 (a+1)Q_{-,a}^*Q_{-,a}
 +aQ_{+,a}^*Q_{+,a}=(2a+1)p(a)I.
 \label{x_reciprocidad:lib02:balance-pesado}
\end{equation}
L'application
\begin{equation}
 W_av=\frac{(\sqrt{a+1}\,Q_{-,a}v,
                    \sqrt a\,Q_{+,a}v)}{\sqrt{(2a+1)p(a)}}
 \label{x_reciprocidad:lib02:W}
\end{equation}
est isométrique et s'inverse sur son image au moyen de \(W_a^*\), de norme un. Pour les canaux non normalisés \(x=Q_{-,a}v\), \(y=Q_{+,a}v\),
\begin{equation}
 v=\frac{x+y}{2a+1},\qquad
 Uv=\frac{ay-(a+1)x}{2a+1}.
 \label{x_reciprocidad:lib02:reconstruccion-bilateral}
\end{equation}
Un couple \((x,y)\) appartient à cette image non normalisée si et seulement si
\(Q_{+,a}x=Q_{-,a}y\).
\end{theorem}
\begin{proof}
Dans \(Q_+-Q_-^*\) apparaît \(I+U+U^2=3P_0\). Développer
\(Q_+Q_-\), \(Q_+^*Q_+\) et utiliser
\(U^2=3P_0-I-U\) donne \eqref{x_reciprocidad:lib02:dos-cartas}. La combinaison
pondérée annule les termes de \(P_0\) dans
\eqref{x_reciprocidad:lib02:gram-Qmenos} et \eqref{x_reciprocidad:lib02:dos-cartas} ; cela prouve
\(W_a^*W_a=I\). Comme \(Q_-+Q_+=(2a+1)I\), la somme des
canaux donne \(v\) ; leur combinaison \(ay-(a+1)x\) donne
\((2a+1)Uv\).

La condition d’image est nécessaire car les deux opérateurs commutent.
Si elle est satisfaite, définissons \(v=(x+y)/(2a+1)\). Alors
\[
 Q_-v=\frac{Q_-x+Q_-y}{2a+1}
 =\frac{(Q_-+Q_+)x}{2a+1}=x,
\]
et de même \(Q_+v=y\). La formule de l'inverse adjoint est celle de toute isométrie ; pour la paire brute, l'inverse linéaire \((x,y)\mapsto(x+y)/(2a+1)\) a pour norme \(\sqrt2/(2a+1)\) dans l'espace produit.
\end{proof}

Dans \(a=8\), les secteurs fixe et d’augmentation ont pour lectures
\((49,100)\) et \((73,73)\), respectivement, et
\[
 9\|Q_-v\|^2+8\|Q_+v\|^2=17\cdot73\|v\|^2=1241\|v\|^2,
 \qquad v=\frac{Q_-v+Q_+v}{17}.
\]
Pour \(s_{+,a}(v)=\|Q_{+,a}v\|^2/\|v\|^2\),
\[
 s_{+,a}=p(a)+3(a+1)w,
 \qquad (a+1)s_{-,a}+a s_{+,a}=(2a+1)p(a).
\]
L’annulation vaut même en présence d’une composante fixe. La composition bilatérale est
une formalisation sur des opérateurs antécédents, et non l’affirmation que
tout retour TPK serait automatiquement cette paire.

Si \(\mathcal V:H\to\mathcal M\) est une mémoire isométrique
complète, \(U^{\mathcal V}=\mathcal VU\mathcal V^*\) satisfait
\((U^{\mathcal V})^3=\Pi_{\mathcal M}\), où
\(\Pi_{\mathcal M}=\mathcal V\mathcal V^*\) est l’identité sur
l’image. Le gramien transporté y est
\[
 (Q_{-,a}^{\mathcal V})^*Q_{-,a}^{\mathcal V}
 =p(a)\Pi_{\mathcal M}-3a\mathcal VP_0\mathcal V^*.
\]
Elle conserve les lectures et la récupération bilatérale. Utiliser des copies diagonales de \(U\) n'est légitime que lorsque les canaux entrelacent cette action ; dans le cas contraire, le transport par \(\mathcal V\) est la définition correcte sur l'image.

\section{Coordonnée de structure fine et itérations entières}
\label{x_reciprocidad:lib02:ensayo-inverso}

Le lecteur complet d’alpha, avec retenue, fournit le préfixe par
blocs
\[
 (007,297,352,569,283,800,997,285,105,472,380,662,999,\ldots).
\]
La fenêtre à frontière nulle qui se termine par 663 est un autre objet ; elle n’est pas
échangée avec la sortie complète. Pour l’essai, l’intervalle
fermé hérité suffit
\begin{equation}
 \frac{7297352569283800997285105472380662}{10^{36}}
 \leq\alpha_A\leq
 \frac{7297352569283800997285105472380663}{10^{36}}.
 \label{x_reciprocidad:lib02:intervalo-alpha}
\end{equation}
La constante n'est identifiée à aucune de ses extrémités rationnelles. Posons \(s_\alpha=10^4\alpha_A\), de sorte que
\begin{equation}
 \begin{gathered}
 \ell\leq s_\alpha\leq h,\\
 \ell:=72.97352569283800997285105472380662,\\
 h:=72.97352569283800997285105472380663.
 \end{gathered}
 \label{x_reciprocidad:lib02:intervalo-s}
\end{equation}
La condition inverse du lecteur fixe est exactement
\begin{equation}
 s(v)=s_\alpha
 \quad\Longleftrightarrow\quad
 w(v)=\frac{73-s_\alpha}{24}
 \simeq0.00110309613175.
 \label{x_reciprocidad:lib02:criterio-alpha}
\end{equation}

Aucun des six états de \cref{x_reciprocidad:lib02:seis-estados} ne satisfait cette
condition. La seconde mémoire a \(s=73\), tandis que les cinq autres
ont \(s<72\). Cette dernière inégalité se vérifie exactement
sous la forme \(24b_v>a_v\). Dans les trois lignes rationnelles, il suffit de multiplier des
entiers ; pour \(u\) et \(q_K\), les différences sont
\[
 24b_u-a_u=\frac{66073183}{20}+\frac{5235904}{5}\sqrt5>0,
\]
\[
 24b_{q_K}-a_{q_K}
 =\frac{45259241}{15}-\frac{5235904}{5}\sqrt5>0.
\]
Pour la seconde, \(\sqrt5<5/2\) donne déjà une borne inférieure positive. L'exclusion concerne ce lecteur et ces états ; elle ne déclare pas une absence globale de relations avec alpha.

Comme \(0<s_\alpha<73\), le rayon
\[
 r_\alpha=\frac{4\sqrt{s_\alpha}}{81}<\frac{4\sqrt{73}}{81}
\]
est une garantie prudente valide, mais ce n’est pas le rayon optimal du
lecteur canonique. Tout nombre entre zéro et 73 donnerait de la même manière
une garantie moindre. Si l’on imposait \(a^2+a+1=s_\alpha\) en changeant
\(a\), \(\sqrt{2a}\) et \((a+1)^2\) changeraient aussi ; on ne
peut pas remplacer seulement le 73 et conserver 4 et 81 comme si l’opérateur
n’avait pas changé.

\begin{theorem}[Itération exacte du poids et récurrence de la lecture]
\label{x_reciprocidad:lib02:iteraciones}
Soit \(v_n=Q_{-,a}^nv_0\), \(a>0\), et \(b=(a-1)^2\). Lorsque \(v_n\ne0\),
\begin{equation}
 w_n=\frac{b^nw_0}{b^nw_0+p(a)^n(1-w_0)}.
 \label{x_reciprocidad:lib02:pesos-iterados}
\end{equation}
Pour \(a\ne1\) et \(0<w_0<1\),
\[
 \frac{w_n}{1-w_n}
 =\left(\frac b{p(a)}\right)^n\frac{w_0}{1-w_0},
 \quad w_n\downarrow0,\quad s_n\uparrow p(a),
\]
et
\begin{equation}
 s_{n+1}=p(a)+b-\frac{bp(a)}{s_n},\qquad
 s_{n+1}-s_n=\frac{(s_n-b)(p(a)-s_n)}{s_n}>0.
 \label{x_reciprocidad:lib02:riccati}
\end{equation}
\end{theorem}
\begin{proof}
L’opérateur et ses adjoints préservent les deux secteurs orthogonaux. Par
\eqref{x_reciprocidad:lib02:gram-Qmenos}, la norme au carré fixe est multipliée par
\(b\) à chaque étape et la complémentaire par \(p(a)\). Cela donne
\eqref{x_reciprocidad:lib02:pesos-iterados} et le rapport des poids. Comme
\(p(a)-b=3a>0\), la monotonie et la limite sont strictes dans le cas
mixte. Enfin,
\(s_n=\|v_{n+1}\|^2/\|v_n\|^2\) ; la suite des normes est
une combinaison de \(b^n\) et \(p(a)^n\), qui vérifie
\(a_{n+2}=(p(a)+b)a_{n+1}-bp(a)a_n\). La division par
\(a_{n+1}\) donne la récurrence et la factorisation de sa différence donne
\eqref{x_reciprocidad:lib02:riccati}.
\end{proof}

Si \(a=1\), le secteur fixe s'annule à la première étape. Lorsque le complément n'est pas nul, la lecture ultérieure vaut exactement trois ; un état purement fixe se transforme en zéro et son quotient cesse d'être défini. Pour \(a\ne1\), l'itération inverse satisfait \(s_{n-1}=bp(a)/(p(a)+b-s_n)\) et diminue la lecture mixte. Inverser l'orientation \(U\mapsto U^{-1}\) laisse, en revanche, inchangée la suite des normes des itérations vers l'avant. Les puissances de \(Q_{+,a}\) multiplient le rapport des poids par \(((a+2)^2/p(a))^n>1\) et diminuent \(s_{-,a}\) dans les états mixtes. Ce sont des filtres réels de l'état, et non des changements de carte ni une trajectoire temporelle TPK identifiée par définition.

\begin{proposition}[Essai complet de toutes les itérations entières]
\label{x_reciprocidad:lib02:cruces-enteros}
Pour \(a=8\), aucune itération entière \(Q_-^nv\), \(n\geq0\),
avec \(v=K,k,m_0\), ne satisfait \(s=s_\alpha\). Les seuls franchissements
de l’intervalle cible se produisent entre les indices \(21,22\),
\(14,15\) et \(15,16\), respectivement. Pour \(v=m_1\), la
lecture demeure exactement 73 dans toutes les itérations.
\end{proposition}
\begin{proof}
Définissons les extrémités rationnelles du poids cible
\[
 w_L=\frac{73-h}{24}
 =\frac{2647430716199002714894527619337}
 {2400000000000000000000000000000000},
\]
\[
 w_H=\frac{73-\ell}{24}
 =\frac{441238452699833785815754603223}
 {400000000000000000000000000000000}.
\]
Pour \(w_0=A/(A+B)\), la formule exacte est
\(w_n=A49^n/(A49^n+B73^n)\). Les données et les croisements sont
\[
\begin{array}{c|r|r|c}
 v&A&B&N\\\hline
 K&10684363&2069408&21\\
 k&3512283&8277632&14\\
 m_0&1052377&1837634&15
\end{array}
\]
et satisfont, dans chaque ligne,
\[
 \frac{A49^N}{A49^N+B73^N}>w_H>w_L>
 \frac{A49^{N+1}}{A49^{N+1}+B73^{N+1}}.
\]
Ce sont des inégalités d’entiers : par exemple, si
\(w_H=p_H/q_H\), la première équivaut à
\((q_H-p_H)A49^N>p_HB73^N\), et la dernière se vérifie
de manière analogue avec \(w_L\). La substitution des trois lignes et des
entiers écrits prouve les franchissements stricts sans arrondi
d’alpha. La monotonie stricte de \cref{x_reciprocidad:lib02:iteraciones} exclut
tous les indices antérieurs et postérieurs ; entre \(N\) et
\(N+1\), il n’y a pas d’autre entier. Pour \(m_1\), \(w_0=0\) demeure
nul par \eqref{x_reciprocidad:lib02:pesos-iterados}.
\end{proof}

Comme référence numérique, les lectures de part et d’autre des croisements sont
\[
\begin{array}{c|c|c}
 v&s_N&s_{N+1}\\\hline
 K&72.97136264842803&72.98077012438719\\
 k&72.96167997982852&72.97426483341773\\
 m_0&72.96527892867391&72.97668298511348
\end{array}
\]
La preuve couvre tous les \(n\geq0\), et non seulement un échantillon fini.
L’orientation inverse donne les mêmes croisements ; les puissances positives de
\(Q_+\) et les puissances inverses de \(Q_-\) diminuent les
lectures initiales, déjà inférieures à 72, et n’atteignent pas non plus la cible
dans ces trois familles. Les produits mixtes \(Q_-^nQ_+^m\)
constituent une autre classe à deux indices, de rapport des poids
\((49/73)^n(100/73)^m w_0/(1-w_0)\) ; ils ne sont pas exclus par
la preuve à un paramètre précédente, et aucun indice n’est choisi ici pour ajuster
alpha.

\section{Réciprocité angulaire et correction régionale existante}
\label{x_reciprocidad:lib02:correccion-regional}

Les changements de carte de \cref{x_reciprocidad:lib02:inversion-carta} incluent la réciprocité et la déformation elliptique déjà construites. Dans un plan réalisé par \(M_\eta\), la forme transportée est \(M_\eta^{-T}M_\eta^{-1}\) ; la similitude ne change pas les valeurs propres du cycle. Maintenir le produit euclidien précédent après le changement de carte définirait un autre observable. L'involution \(q_+\leftrightarrow q_-\) conserve \(A\), inverse \(C^*\) et envoie \(\eta_{\rm el}\) sur \(-\eta_{\rm el}\). La relation polaire de l'ellipse \(\tan\phi=e^{-\eta_{\rm el}}\tan\theta\) change le paramétrage de la courbe, et non la phase spectrale d'un transport semblable. Elle n'ajoute pas de petite phase à \(S^4\).

La conjugaison incidentielle a un autre type. Sa loi de blocs conserve
\[
 z_m^\dagger=z_{\rho(m)}
 +100\mathbf1_{\{c_{\rho(m)}\ne0\}}-20c_{\rho(m)},
 \qquad c_m=[C_m]_{10},
\]
avec les termes de frontière avant la réduction décimale. Ce n’est pas en
général une isométrie du vecteur de blocs. L’appliquer au registre
requiert son raccordement au registre des incidences et à ses extrémités ; on ne la
remplace pas par une correction angulaire inventée.

\begin{proposition}[Correction régionale du contre-angle]
\label{x_reciprocidad:lib02:retorno-regional}
Les coordonnées antécédentes d’action et d’angle vérifient
\[
 A_{\rm deg}=1000\alpha_A,\qquad
 C^*_{\rm ret,deg}=2(\eta_{\rm ret}+\alpha_A),\qquad
 C^*_{\rm pre,deg}=2(H_5+\alpha_A),
\]
\[
 \eta_{\rm ret}=H_5-\frac{90}{\pi}\mathcal C_\pi(\alpha_A),
 \quad
 \mathcal C_\pi(\alpha_A)=169\alpha_A^6+
 \frac{D_A\alpha_A^7}{1-D_A\alpha_A},\qquad
 D_A=\exp(-100\pi\alpha_A/9).
\]
La différence en radians est exactement
\[
 \theta_{C,\rm pre}-\theta_{C,\rm ret}=\mathcal C_\pi(\alpha_A).
\]
\end{proposition}
\begin{proof}
La différence en degrés est \(2(H_5-\eta_{\rm ret})=180\mathcal C_\pi/\pi\). Multiplier une seule fois par \(\pi/180\) donne le résultat.
\end{proof}

Ici, \(\pi\) et \(\alpha_A\) sont les coordonnées déjà engendrées ;
\(H_5\) et les autres lecteurs conservent leurs définitions
antécédentes, et non un recalibrage. Les évaluations ultérieures sont
\[
\begin{array}{c|c}
 \text{coordonnée}&\text{radians, environ}\\\hline
 \theta_A&0.127362829013\\
 \theta_C&0.0370662264658\\
 \theta_C/6\text{ dans chaque paire antipodale}&0.00617770441097\\
 \mathcal C_\pi&2.55207421805\cdot10^{-11}
\end{array}
\]
Les canaux \(q_\pm\) sont des atténuations réelles positives ; leur exposant
ne devient pas automatiquement une phase unitaire au moyen d’un facteur
\(i\). On n’introduit pas de diviseurs supplémentaires pour approcher les
échelles de l’essai qui suit.

\section{Problème inverse angulaire dans le registre dodécaphasique}
\label{x_reciprocidad:lib02:ensayo-angular}

Pour une cible générale \(s\in[49,81]\), l'orientation principale du secteur pur s'obtient à partir de
\[
 |8-e^{i\theta}|^2=65-16\cos\theta,
 \qquad \theta(s)=\arccos\frac{65-s}{16}\in[0,\pi].
\]
En particulier, la donnée ultérieure \(s_\alpha\) détermine
\[
 \theta_\alpha=\theta(s_\alpha),\qquad
 \delta_\alpha=\frac{2\pi}{3}-\theta_\alpha.
\]
Cette formule résout un problème inverse avec une sortie déjà produite ;
elle ne prétend pas produire alpha de nouveau.

\begin{lemma}[Structure angulaire du complément ternaire]
\label{x_reciprocidad:lib02:J}
L'opérateur \(J=(U-U^2)/\sqrt3\) satisfait
\[
 J^*=-J,\quad J^2=-P,\quad JP_0=P_0J=0,\quad JP=PJ=J,
 \qquad U=P_0-\frac12P+\frac{\sqrt3}{2}J.
\]
\end{lemma}
\begin{proof}
L’adjoint échange \(U,U^2\). En outre,
\((U-U^2)^2=U+U^2-2I=-3P\), et
\((I+U+U^2)(U-U^2)=0\). Les autres relations se déduisent de
\(P=I-P_0\) ; additionner \((U+U^2)/2\) et \((U-U^2)/2\)
donne la dernière expression.
\end{proof}

\begin{theorem}[Famille angulaire et propriétés conservées]
\label{x_reciprocidad:lib02:Utheta}
Pour \(\theta\in\mathbb R\), définissons
\[
 U_\theta=P_0+\cos\theta\,P+\sin\theta\,J,\qquad
 Q_{-,\theta}=8I-U_\theta.
\]
On a
\[
 U_\theta^*U_\theta=I,\qquad
 U_\theta U_\varphi=U_{\theta+\varphi},\qquad U_{2\pi/3}=U,
\]
\begin{equation}
 Q_{-,\theta}^*Q_{-,\theta}
 =49P_0+(65-16\cos\theta)P.
 \label{x_reciprocidad:lib02:gram-angular}
\end{equation}
Les lecteurs \(T_\theta=(8I+U_\theta)/9\), \(D_\theta=\sqrt8(U_\theta-I)/9\) conservent \(T_\theta^*T_\theta+D_\theta^*D_\theta=I\). En revanche,
\[
 U_\theta^3=I\ \Longleftrightarrow\ 3\theta\in2\pi\mathbb Z,
 \qquad
 \det Q_{-,\theta}=7^4(65-16\cos\theta)^4.
\]
\end{theorem}
\begin{proof}
Les relations de \cref{x_reciprocidad:lib02:J} annulent les termes entre secteurs.
Dans le complément, \(J^2=-I\), de sorte que
\(U_\theta^*U_\theta=P_0+(\cos^2\theta+\sin^2\theta)P=I\).
Multiplier deux expressions et utiliser les formules d’addition donne la
loi de composition. L’expression de \(U\) dans le lemme donne la valeur
ternaire. Puisque
\(U_\theta+U_\theta^*=2P_0+2\cos\theta P\), développer le
Gram donne \eqref{x_reciprocidad:lib02:gram-angular}. Les développements
\[
 81T_\theta^*T_\theta=65I+8(U_\theta+U_\theta^*),\qquad
 81D_\theta^*D_\theta=16I-8(U_\theta+U_\theta^*)
\]
prouvent le bilan. La loi de composition donne \(U_\theta^3=P_0+\cos(3\theta)P+\sin(3\theta)J\), qui est l'identité exactement sous la condition indiquée.

Pour le déterminant, \(Q_{-,\theta}\) agit par 7 sur les quatre
dimensions fixes. Dans le complément, un vecteur unitaire \(e\) et
\(Je\) forment un couple orthonormé ; leur complément orthogonal est stable
sous \(J\). En répétant, on obtient quatre plans, sur chacun
desquels le déterminant est
\((8-\cos\theta)^2+\sin^2\theta=65-16\cos\theta\).
Multiplier les contributions démontre la formule.
\end{proof}

\begin{corollary}[Solution inverse et séparation des corrections]
\label{x_reciprocidad:lib02:fase-alpha}
Dans \(\operatorname{im}P\), le lecteur \(Q_{-,\theta_\alpha}\) a pour rapport quadratique exactement \(s_\alpha\). Les orientations \(\pm\theta_\alpha\) donnent le même rapport et
\[
 73-s_\alpha
 =8(1-\cos\delta_\alpha+\sqrt3\sin\delta_\alpha).
\]
On a \(U_{\theta_\alpha}^3\ne I\), et \(\delta_\alpha\ne\mathcal C_\pi(\alpha_A)\).
\end{corollary}
\begin{proof}
La définition de la phase résout le coefficient complémentaire de
\eqref{x_reciprocidad:lib02:gram-angular} ; le cosinus est pair. La substitution de
\(\theta_\alpha=2\pi/3-\delta_\alpha\) donne l’identité stable
pour la différence. Comme \(72<s_\alpha<73\),
\(\pi/2<\theta_\alpha<2\pi/3\) et
\(0<3\delta_\alpha<\pi/2\) ; par conséquent, \(3\theta_\alpha\)
n’est pas un multiple de \(2\pi\).

La séparation des corrections peut être prouvée sans se fier à des chiffres approchés. De l'intervalle résulte \(73-s_\alpha>0.026\). Si \(\delta_\alpha\leq10^{-3}\), les inégalités \(1-\cos t\leq t^2/2\), \(\sin t\leq t\) et \(\sqrt3<2\) donneraient
\[
 73-s_\alpha\leq8(10^{-6}/2+2\cdot10^{-3})=0.016004,
\]
contradiction. Ainsi, \(\delta_\alpha>10^{-3}\). D’autre part,
\(0<\alpha_A<10^{-2}\), \(0<D_A<1\), impliquent
\[
 0<\mathcal C_\pi(\alpha_A)
 <169\cdot10^{-12}+\frac{10^{-14}}{1-10^{-2}}
 <170\cdot10^{-12}<10^{-7}.
\]
Les deux quantités sont distinctes.
\end{proof}

L'évaluation diagnostique donne
\[
 \delta_\alpha\simeq0.00190956706826\ {\rm rad}
 \simeq0.109410133709^\circ,\qquad
 \theta_\alpha\simeq119.890589866291^\circ.
\]
La correction régionale est environ \(7.48\cdot10^7\) fois
plus petite. La famille inverse conserve l’orthogonalité et le bilan, mais perd
l’ordre trois à cette phase et change le déterminant original. Sa carte
réelle ne fournit pas à elle seule une matrice entière ni le quotient réticulaire
d’indice 73. Le fait de résoudre la donnée prescrite ne démontre pas
\(73=10^4\alpha_A\).

\begin{proposition}[Atteignabilité en conservant l'état]
\label{x_reciprocidad:lib02:alcanzabilidad}
Pour \(v\ne0\), avec \(w=w(v)\) fixé,
\[
 s_\theta(v)=49w+(65-16\cos\theta)(1-w).
\]
Si \(w<1\), son image est exactement
\([49,81-32w]\). Il existe donc une phase qui atteint
\(s_\alpha\) si et seulement si
\[
 w\leq\frac{81-s_\alpha}{32}\simeq0.250827322099.
\]
Lorsqu'elle existe, son orientation principale est
\[
 \theta_v=\arccos\left[
 \frac1{16}\left(65-\frac{s_\alpha-49w}{1-w}\right)\right].
\]
Si \(w=1\), la lecture est toujours 49 et n'atteint pas la cible.
\end{proposition}
\begin{proof}
Évaluer \eqref{x_reciprocidad:lib02:gram-angular} et diviser par \(\|v\|^2\)
donne la formule. Le cosinus parcourt \([-1,1]\), de sorte que le
coefficient complémentaire parcourt \([49,81]\). Sa combinaison avec
le poids fixe donne l’intervalle complet, sans trous. Résoudre l’égalité
et vérifier le domaine de l’arccosinus donne la condition et la phase.
Le cas \(w=1\) s’évalue sans division par \(1-w\).
\end{proof}

\begin{corollary}[Quatre impossibilités et deux phases possibles]
\label{x_reciprocidad:lib02:seis-angulares}
En maintenant l'état et le secteur fixe, la famille angulaire n'atteint pas \(s_\alpha\) sur \(K,k,m_0,u\). Elle l'atteint en revanche sur \(m_1\) et \(q_K\).
\end{corollary}
\begin{proof}
Du tableau exact \eqref{x_reciprocidad:lib02:tabla-normas}, on obtient
\[
\begin{array}{c|c}
 v&32b_v-9a_v\\\hline
 K&227115677/3\\
 k&2094607/4\\
 m_0&122655440/243\\
 u&37201759/20+(7859008/15)\sqrt5
\end{array}
\]
Toutes les entrées sont positives. Ainsi \(w>9/32\), et le maximum \(81-32w\) est inférieur à \(72<s_\alpha\). Pour \(m_1\), \(b_{m_1}=0\) et \(a_{m_1}>0\), donc l'intervalle est \([49,81]\). Pour \(q_K\),
\[
 a_{q_K}-4b_{q_K}=\frac{588839+1195712\sqrt5}{15}>0.
\]
Par conséquent, \(w(q_K)<1/4\) et son maximum est supérieur à
\(73>s_\alpha\). Dans les deux cas, la cible appartient à
l’intervalle atteignable.
\end{proof}

Les évaluations des maxima et, lorsqu’elles existent, les phases sont
\[
\begin{array}{c|c|c}
 v&\max_\theta s_\theta(v)&\theta_v\\\hline
 K&54.1922726227&\text{n'existe pas}\\
 k&71.4670172771&\text{n'existe pas}\\
 m_0&69.3474270513&\text{n'existe pas}\\
 u&66.8296175021&\text{n'existe pas}\\
 m_1&81&119.8905898663^\circ\\
 q_K&77.3923674871&133.5296853010^\circ
\end{array}
\]
La phase du secteur pur ne convient pas indistinctement à un état
mixte : sa lecture serait
\(49w+s_\alpha(1-w)=s_\alpha-(s_\alpha-49)w\).
Les lecteurs originaux de \(V_2\) sont des polynômes en \(S\), comme
\(P_0,P,J,U_\theta\). Ils entrelacent donc également la famille
angulaire, et \(s_\theta(V_2v)=s_\theta(v)\). Archiver les
compléments conserve la lecture conjointe, mais n’élimine pas les
restrictions d’atteignabilité.

\section{Composition des lecteurs et portée des problèmes inverses}
\label{x_reciprocidad:lib02:alcance}

Le facteur \(73\) a été déterminé par trois opérations
compatibles et différenciées : le gramien de deux lecteurs nonadiques, la
similitude du module d’augmentation ternaire et les déterminants entiers des
deux cartes conjuguées. Le bilan bilatéral conserve l’état
complet et admet une inversion explicite ; la mémoire complète conserve les
lectures transportées, tandis que la sélection d’un canal peut changer
ses poids. Les rayons de correction sont uniformes dans tous les registres
entiers du domaine algébrique et dans toute la famille \(a>0\), avec
des preuves de coupes et de parité qui ne dépendent pas d’un échantillonnage.

L'essai inverse ajoute des conditions vérifiables, et non un calibrage caché. L'inversion de l'orientation et la covariance métrique ne changent pas la lecture. Les itérations entières spécifiées sont résolues pour tous leurs indices par un croisement exact et la monotonie. La variation angulaire produit une famille qui réalise la cible dans les états permis, mais la donnée \(s_\alpha\) sélectionne sa phase et l'indice ternaire original est perdu. Son identification à une réalisation antérieure de \(\Gamma_9\) exigerait une composition typée sur l'état enrichi ; cet essai n'apporte pas une telle identification et ne déclare pas son absence dans tout le corpus.

Les nouvelles formalisations réunissent des opérateurs et des questions antécédents ;
aucune priorité mondiale n’est attribuée au codage, à l’utilisation de
vecteurs cycliques ou à la décomposition polaire. Les chiffres décimaux
exposés sont des évaluations ultérieures de formules exactes. Les
conditions de non-nullité, de charge, de métrique, de réseau et de normalisation font partie
du résultat, et leurs contre-exemples restent aux côtés des preuves.
Ainsi est conservée la fonction de la constante holographique de couplage
arithmético-géométrique : un objet produit qui articule lecteurs, mémoire
et incidence, sans transformer chaque coïncidence numérique en une nouvelle loi
physique.
